  4626	\subsection{Propriétés d'exactitude de $\Ind(C)$}\etiquette[8.9]{I.sec8.9}
  4627	
  4628	\begin{enonce*}{Proposition 8.9.1}\etiquette[8.9.1]{I.prop8.9.1}
  4629	Soit $C$ une $𝒰$-catégorie.
  4630	\begin{enumerate}
  4631	\renewcommand{\theenumi}{\alph{enumi}}
  4632	\renewcommand{\labelenumi}{\theenumi)}
  4633	\item Les foncteurs canoniques $L$ (8.2.4.8) et $c$ (8.4.1)
  4634	$$
  4635	C\xrightarrow{c}\Ind(C)\xrightarrow{L}\widehat{C}
  4636	$$
  4637	commutent aux limites projectives.
  4638	
  4639	\item Supposons que dans $C$ les limites inductives finies soient
  4640	représentables, et que $C$ soit équivalente à une petite
  4641	catégorie.
  4642	Alors dans $\Ind(C)$ les petites limites projectives sont
  4643	représentables.
  4644	
  4645	\item Si\pageoriginale dans $C$ les petites limites projectives
  4646	(\resp les limites projectives finies) sont représentables, alors il
  4647	en est de même dans $\Ind(C)$. Soit $J$ une catégorie qui est finie
  4648	et rigide, ou discrète; si les limites projectives de type $J$ sont
  4649	représentables dans $C$, ce même type de limites projectives est
  4650	représentables dans $\Ind(C)$.
  4651	
  4652	\item Les petites limites inductives filtrantes dans $\Ind(C)$ (elles
  4653	sont représentables en vertu de \hyperref[I.prop8.5.1]{8.5.1}) sont exacts à
  4654	gauche, \ie commutent aux limites projectives finies.
  4655	\end{enumerate}
  4656	\end{enonce*}
  4657	
  4658	\noindent \emph{Démonstration}.
  4659	\begin{enumerate}
  4660	\renewcommand{\theenumi}{\alph{enumi}}
  4661	\renewcommand{\labelenumi}{\theenumi)}
  4662	\item Le fait que $L$ commute aux limites projectives résulte du
  4663	fait qu'il est pleinement fidèle et du calcul des limites
  4664	projectives dans $\widehat{C}$ « argument par argument », qui
  4665	implique que pour vérifier qu'un système projectif de morphismes
  4666	$F\rightarrow F_i$ $(i\in I)$ dans $\widehat{C}$ fait de $F$ une
  4667	limite projective des $F_i$, il suffit de vérifier que l'assertion
  4668	analogue est vraie pour les %%
  4669	systèmes projectifs d'applications ensemblistes
  4670	$\Hom(X,F)\rightarrow \Hom(X,F_i)$, pour tout $X\in \Ob C$. Or
  4671	$C\subset \Ind(C)$.
  4672	
  4673	Le fait que $c$ commute aux limites projectives résulte formellement
  4674	du fait que $L$ y commute, ainsi que le composé $L\circ
  4675	c:C\rightarrow \widehat{C}$.
  4676	
  4677	\item Compte tenu de a), l'assertion revient à dire que toute
  4678	limite projective de préfaisceaux ind-représentables est
  4679	ind-représentable, ce qui résulte aussitôt du critère
  4680	\hyperref[I.them8.3.3]{8.3.3} (v),
  4681	compte tenu qu'une limite projective de foncteurs exacts à gauche
  4682	est exact à gauche (ce qui résulte du fait que « les limites
  4683	projectives commutent aux limites projectives » (\hyperref[I.sec2.5.0]{2.5.0})).
  4684	
  4685	\item Résulte\pageoriginale formellement de la propriété
  4686	  analogue de
  4687	$\widehat{C}$ (\hyperref[I.cor3.3]{3.3}), et du fait que $L$ est conservatif (étant
  4688	pleinement fidèle) et commute aux limites du type envisagé (a) et
  4689	8.5.1).
  4690	
  4691	\item Il est bien connu (\cf \hyperref[I.prop2.3]{2.3}) que la
  4692	  représentabilité des
  4693	limites projectives finies équivaut à celle des limites projectives
  4694	des types
  4695	
  4696	\begin{center}
  4697	vide $, \ldots ,\quad \vcenter{ \xymatrix@C=.25cm@R=.5cm{
  4698	                     \bigdot \ar[rd]&       &\bigdot\ar[ld]\\
  4699	                             &\bigdot&
  4700	}}$
  4701	\end{center}
  4702	
  4703	\noindent (correspondants à des ensembles ordonnés finis
  4704	particuliers), et celle des petites limites projectives revient à
  4705	celle des produits et des limites projectives finies. Donc la
  4706	deuxième assertion faite dans (c) implique la première.
  4707	\end{enumerate}
  4708	
  4709	\emph{Supposons d'abord $I$ fini.} Utilisant le résultat
  4710	\hyperref[I.prop8.8.5]{8.8.5} sur la représentabilité des %%
  4711	foncteurs $\varphi:J\rightarrow \Ind(C)$ sous forme indicielle
  4712	(\hyperref[I.sec8.5.4]{8.5.4})
  4713	$$
  4714	\displaylines{\rlap{$(8.9.1.1)$}\hfill \phi:J\times I\longrightarrow
  4715	C{\quoi},\hfill}
  4716	$$
  4717	l'existence des $\varprojlim \varphi$ est donc un cas particulier du
  4718	résultat plus précis et plus général:
  4719	
  4720	\begin{enonce*}{Corollaire \sisi{8.8.2}{8.9.2}}\etiquette[8.9.2]{I.cor8.9.2}
  4721	Considérons un foncteur $\varphi:J\rightarrow \Ind(C)$ donné sous
  4722	forme indicielle $\Phi$ (8.9.1.1), $J$ étant une
  4723	catégorie finie.
  4724	Si pour tout $i\in \Ob I$, le foncteur partiel $j\rightarrow
  4725	\Phi(j,i)$ a une limite projective (\resp inductive) représentable
  4726	dans $C$, alors $\varphi$ a une limite projective (\resp inductive)
  4727	représentable dans $\Ind(C)$, et on a un isomorphisme canonique
  4728	$$
  4729	\displaylines{\rlap{$(8.9.2.1)$}\hfill
  4730	\varprojlim\varphi\simeq\textrm{« }{\varinjlim_i}\textrm{ »} \varprojlim_j
  4731	\Phi(j,i)\hfill}
  4732	$$
  4733	(\resp\pageoriginale
  4734	$$
  4735	\displaylines{\rlap{$(8.9.2.2)$} \hfill \varinjlim \varphi \simeq
  4736	\textrm{« }{\varinjlim_i}\textrm{ »} \varprojlim_j\Phi(j,i)){\quoi},\hfill}
  4737	$$
  4738	où $\varprojlim_j$ (\resp. $\varinjlim_j$) est calculé dans $C$.
  4739	\end{enonce*}
  4740	
  4741	Pour la première formule, on utilise simplement que dans $\Ind(C)$
  4742	les limites inductives filtrantes commutent à $\varprojlim_j$, et
  4743	que $c$ commute à $\varprojlim_J$ (8.9.1 d) et a)):
  4744	$$
  4745	\varprojlim \varphi=\varprojlim_j \textrm{« }{\varinjlim_i}\textrm{ »\,} \Phi
  4746	(j,i) \xleftarrow{\sim} \textrm{« }{\varinjlim_i}\textrm{ »\,}
  4747	{\varprojlim_j}^{\Ind(C)} \Phi(j,i)\simeq \textrm{« }{\varinjlim_i}\textrm{ »}%%
  4748	\varprojlim_j{^C \ \Phi(j,i){\quoi}}.
  4749	$$
  4750	La seconde se prouve de même, en utilisant la commutation des
  4751	limites inductives de type $I$ aux limites inductives de type $J$
  4752	(\hyperref[I.sec2.5.0]{2.5.0}) et le fait que $c$ commute aux limites inductives de type
  4753	$J$, prouvé dans \hyperref[I.rem8.9.4]{8.9.4} a) ci-dessous.
  4754	
  4755	\emph{Cas $J$ discret}. Ce cas est contenu dans l'assertion plus
  4756	précise suivante:
  4757	
  4758	\begin{enonce*}{Corollaire 8.9.3}\etiquette[8.9.3]{I.cor8.9.3}
  4759	Soient $I_\alpha$ des catégories filtrantes essentiellement petites
  4760	indexées par un petit ensemble $A$, et pour tout $\alpha \in A$,
  4761	soit $X(\alpha)=(X(\alpha)_{i_\alpha})_{i_\alpha \in I_\alpha}$ un
  4762	ind-objet de $C$ indexé par $I_\alpha$. Soit $I$ la catégorie
  4763	produit des $I_\alpha$, et supposons que pour tout
  4764	$\sheaf{i}=(i_\alpha)_{\alpha\in A}\in I$, le produit
  4765	$\prod_{\alpha\in A}%%
  4766	X (\alpha)_{i_\alpha}$ soit représentable dans $C$. Alors le produit
  4767	$\prod_{\alpha \in A} %%
  4768	X (\alpha)$ est représentable dans $\Ind(C)$, et il est
  4769	canoniquement isomorphe à l'ind-objet indexé par $I$ donné par la
  4770	formule
  4771	$$
  4772	\displaylines{\rlap{$(8.9.3.1)$}\hfill \prod_{\alpha \in A}\;
  4773	\mathop{\textrm{« }\varinjlim\textrm{ »}}_{i_\alpha\in I_\alpha}\;
  4774	X(\alpha)_{i_\alpha}\simeq
  4775	\mathop{\textrm{« }\varinjlim\textrm{ »}}_{\sheaf{i}=(i_\alpha)_{\alpha\in
  4776	A}\in I}\; \left(\prod_{\alpha\in
  4777	A}A(\alpha)_{i_\alpha}\right){\quoi},\hfill}
  4778	$$
  4779	où\pageoriginale le produit du deuxième membre est le produit
  4780	calculé dans $C$. (On notera que $I$ est filtrant et essentiellement
  4781	petit, les $I_\alpha$ l'étant.)
  4782	\end{enonce*}
  4783	
  4784	En effet, il et bien connu (et immédiat par réduction au cas où on
  4785	travaille dans la catégorie des ensembles) que la formule envisagée
  4786	est valable quand on calcule les limites dans $\widehat{C}$. La
  4787	conclusion résulte alors du fait que $L$ commute aux limites
  4788	envisagées (\hyperref[I.prop8.9.1]{8.9.1} a) et \hyperref[I.prop8.5.1]{8.5.1}).
  4789	
  4790	\begin{enonce*}[remark]{Remarques 8.9.4}\etiquette[8.9.4]{I.rem8.9.4}
  4791	La démonstration donnée de 8.9.2, 8.9.3 montre, plus
  4792	généralement,
  4793	que si pour tout $i\in \Ob I$, $\varprojlim_j \Phi(j,i)$ (\resp
  4794	$\varinjlim_j \Phi(j,i)$) \emph{calculé dans} $\Ind(C)$ est
  4795	représentable, alors il en est de même de $\varprojlim_i\varphi$
  4796	(\resp de $\varinjlim_j\varphi$). Ceci et l'argument de c) montre
  4797	que pour une catégorie %%
  4798	donnée $J$ provenant d'un ensemble ordonné fini ou discret (ou plus
  4799	généralement, qui est $C$-admissible (\hyperref[I.prop8.3.1]{8.3.1})), les
  4800	limites projectives (\resp inductives) de type $J$ sont
  4801	représentables dans $\Ind(C)$ si (et seulement si) pour tout
  4802	foncteur $\varphi:J\rightarrow C$, la limite projective (\resp
  4803	inductive) de $\varphi$ \emph{calculée dans} $\Ind(C)$ est
  4804	représentable. Dans le cas non respé, cela signifie aussi, en vertu
  4805	de (a), que toute limite projective de type $J$ de préfaisceaux
  4806	représentables sur $C$ est ind-représentable.
  4807	\end{enonce*}
  4808	
  4809	\begin{enonce*}{Proposition 8.9.5}\etiquette[8.9.5]{I.prop8.9.5}
  4810	Soit $C$ une $𝒰$-catégorie.
  4811	\begin{enumerate}
  4812	\renewcommand{\theenumi}{\alph{enumi}}
  4813	\renewcommand{\labelenumi}{\theenumi)}
  4814	\item Le foncteur canonique
  4815	$$
  4816	c:C\longrightarrow \Ind(C)
  4817	$$
  4818	est exact à droite (donc exact, compte tenu de $8.9.1$ {\rm a)}).
  4819	
  4820	\item Si les limites inductives finies (\resp les sommes finies) sont
  4821	représentables dans $C$, alors les petites limites inductives (\resp
  4822	les petites sommes) sont représentables dans $\Ind(C)$. Soit $J$ un
  4823	ensemble préordonné\pageoriginale fini ou discret; si les limites
  4824	inductives de type $J$ sont représentables dans $C$, il en est de
  4825	même dans $\Ind(C)$.
  4826	\end{enumerate}
  4827	\end{enonce*}
  4828	
  4829	\noindent \emph{Démonstration.}
  4830	\begin{enumerate}
  4831	\renewcommand{\theenumi}{\alph{enumi}}
  4832	\renewcommand{\labelenumi}{\theenumi)}
  4833	\item Supposons que dans $C$ on ait $X\simeq \varinjlim_J X_j$, où
  4834	$J$ est une catégorie finie, $X$ et les $X_j$ dans $C$. On a alors,
  4835	pour tout ind-objet $𝒴=(Y_i)_{i\in I}$ de $C$:
  4836	\begin{gather*}
  4837	\Hom(X,𝒴)\simeq \varinjlim_i \Hom(X,Y_i)\simeq \varinjlim_i
  4838	\varprojlim_j \Hom(X_j,Y_i)\\
  4839	\mathrel{\mathop\simeq^{2.8}}\varprojlim_j \varinjlim_i \Hom(X_j,
  4840	Y_i) \simeq \varprojlim_j \Hom (X_j,𝒴){\quoi},
  4841	\end{gather*}
  4842	et la conclusion voulue résulte de la comparaison des termes
  4843	extrêmes.
  4844	
  4845	\item On a déjà noté dans \hyperref[I.rem8.8.4]{8.8.4} que les
  4846	  assertions faites
  4847	résultent de a) et de \hyperref[I.cor8.9.2]{8.9.2}, du moins pour les limites
  4848	inductives finies. Pour prouver les conclusions faites dans les cas
  4849	infinis, on est ramené au cas des sommes, qui peut s'interpréter
  4850	comme une limite inductive filtrante de sommes finies, et on conclut
  4851	donc grâce à \hyperref[I.prop8.5.1]{8.5.1}.
  4852	\end{enumerate}
  4853	
  4854	\begin{enonce*}[remark]{Remarque 8.9.6}\etiquette[8.9.6]{I.rem8.9.6}
  4855	On a déjà observé que le foncteur $c$ ne commute pas en général aux
  4856	limites inductives filtrantes, donc pas non plus aux sommes
  4857	infinies, contrairement à ce qui a lieu pour $L:\Ind(C)\rightarrow
  4858	\widehat{C}$. Par contre, à l'exception de la commutation aux
  4859	limites inductives \emph{filtrantes} (\hyperref[I.prop8.5.1]{8.5.1}), $L$ n'a
  4860	pratiquement jamais de propriétés de commutation à quelque autre
  4861	type de limites inductives (objet initial, somme de deux objets
  4862	conoyaux de doubles flèches). Ainsi, si $C$ admet un objet initial
  4863	$∅_C$, qui est donc objet initial de $\Ind(C)$ en vertu de
  4864	a), $∅_C$ n'est jamais objet initial de $\widehat{C}$ (\ie
  4865	identique au préfaisceau constant de valeur $∅$), puisque
  4866	$\Hom(∅_C,∅_C)\neq ∅$!
  4867	\end{enonce*}
  4868	
  4869	\begin{enonce*}{Proposition 8.9.7}\etiquette[8.9.7]{I.prop8.9.7}
  4870	Soit\pageoriginale
  4871	$$
  4872	f:C\longrightarrow C'
  4873	$$
  4874	un foncteur entre $𝒰$-catégories, et soit $J$ une catégorie
  4875	finie et rigide $(8.8.5)$. Si les limites inductives (\resp
  4876	projectives) de type $J$ sont représentables dans $C$ et si $f$ y
  4877	commute, alors $\Ind(f):\Ind(C)\rightarrow \Ind(C')$ commute
  4878	également à ce type de limites.
  4879	\end{enonce*}
  4880	
  4881	Cela résulte aussitôt de \hyperref[I.prop8.8.5]{8.8.5} et du calcul
  4882	\hyperref[I.cor8.9.2]{8.9.2} des limites finies dans une catégorie de
  4883	ind-objets, pour un foncteur représenté sous forme indicielle.
  4884	
  4885	\begin{enonce*}{Corollaire 8.9.8}\etiquette[8.9.8]{I.cor8.9.8}
  4886	Si dans $C$ les limites inductives (\resp projectives) finies sont
  4887	représentables, et si $f$ est exact à droite (\resp à gauche) alors
  4888	il en est de même de $\Ind(f)$; dans le cas non respé, $\Ind(f)$
  4889	commute même aux petites limites inductives quelconques.
  4890	\end{enonce*}
  4891	
  4892	La première assertion résulte de \hyperref[I.prop8.9.7]{8.9.7}. La deuxième
  4893	résulte de la première, compte tenu que $\Ind(f)$ commute aux
  4894	limites inductives filtrantes (\hyperref[I.prop8.6.3]{8.6.3}).
  4895	
  4896	\begin{enonce*}[remark]{Exercice 8.9.9}\etiquette[8.9.9]{I.exer8.9.9}
  4897	Soit $C$ une $𝒰$-catégorie.
  4898	\begin{enumerate}
  4899	\renewcommand{\theenumi}{\alph{enumi}}
  4900	\renewcommand{\labelenumi}{\theenumi)}
  4901	\item Supposons que dans $C$ les sommes finies sont
  4902	représentables. Montrer que si dans $C$ les sommes \emph{finies}
  4903	sont disjointes (\resp universelles) (\cf II 4.5), alors dans
  4904	$\Ind(C)$ les \emph{petites} sommes sont disjointes (\resp
  4905	universelles).
  4906	
  4907	\item Supposons que dans $C$ tout morphisme se factorise en un
  4908	épimorphisme suivi d'un monomorphisme (\resp en un épimorphisme
  4909	effectif suivi\pageoriginale d'un monomorphisme, \resp en un
  4910	épimorphisme suivi d'un monomorphisme effectif, \resp un
  4911	épimorphisme effectif suivi d'un monomorphisme effectif). Montrer
  4912	que $\Ind(C)$ satisfait à la même propriété. Dans le dernier cas
  4913	respé, on conclut donc que tout épimorphisme de $\Ind(C)$ est
  4914	effectif, tout monomorphisme de $\Ind(C)$ est effectif, tout
  4915	bimorphisme de $\Ind(C)$ est un isomorphisme; montrer que dans ce
  4916	cas tout épimorphisme (\resp monomorphisme) de $\Ind(C)$ peut se
  4917	représenter par un système inductif d'épimorphismes (\resp de
  4918	monomorphismes) de $C$. Si on suppose de plus que dans $C$ tout
  4919	épimorphisme (\resp tout monomorphisme) est \emph{universel}, la
  4920	même propriété est vraie dans $\Ind(C)$.
  4921	
  4922	\item Supposons $C$ additive (\resp abélienne), alors $\Ind(C)$
  4923	l'est également.
  4924	
  4925	\item Soient $X=\textrm{« }\varinjlim_I\textrm{ »\,}X_i$ un ind-objet de $C$,
  4926	$R\rightrightarrows X$ une relation d'équivalence dans $X$. Pour
  4927	tout $i\in \Ob I$, soit $R_i\rightrightarrows X_i$ la relation
  4928	d'équivalence dans $X_i$ (considéré comme objet de $\Ind(C)$)
  4929	induite par $R$ via $X_i\rightarrow X$. (On suppose que le produit
  4930	fibré $R_i=R\times_{X_XX}X_i \times X_i$ dans $(\Ind(C))$ est
  4931	représentable par un objet de $\Ind(C)$, ce qui est le cas si dans
  4932	$C$ les limites projectives finies sont représentables.) Montrer que
  4933	pour que $R$ soit effective, il suffit que les $R_i$ le soient.
  4934	
  4935	\item Soient $X$ un objet de $C$, $R$ une relation d'équivalence
  4936	dans $c(X)$ \ie dans $X$ considéré comme objet de $\Ind(C)$.
  4937	Supposons que dans $C$ les produits fibrés soient représentables.
  4938	Pour que $R$ soit effective, il faut que $R$ soit de la forme
  4939	$\textrm{« }\varinjlim_i\textrm{ »\,} R_i$, où les $R_i$ sont des relations
  4940	d'équivalence dans $X$ (regardé comme un objet de $C$), et cette
  4941	condition est\pageoriginale suffisante si on suppose que dans $C$
  4942	les relations d'équivalence sont effectives.
  4943	
  4944	\item Supposons que dans $C$ tout morphisme se factorise en un
  4945	épimorphisme effectif suivi d'un monomorphisme effectif, et que les
  4946	limites projectives finies soient représentables. Soit $X$ un objet
  4947	de $C$, et $R$ un sous-objet de $X\times X$, regardé comme élément
  4948	de $\Ind(C)$. Écrivons $R$ sous la forme $\textrm{« }\varinjlim\mbox{''
  4949	}Z_i$, où $(Z_i)_{i\in I}$ est une famille filtrante croissante de
  4950	sous-objets de $X\times X$ \emph{dans} $C$ (c'est possible grâce a
  4951	b)). Montrer que pour que $R$ soit une relation d'équivalence dans
  4952	$X$, il faut et il suffit que pour tout $i\in \Ob I$, existe un
  4953	$j\in \Ob I$ qui contienne $s(Z_i)$ et $Z_i\circ Z_i$, où $s$ est la
  4954	symétrie de $X\times X$.
  4955	
  4956	\item Supposons que dans $C$ les limites projectives finies ainsi que
  4957	les sommes finies soient représentables, que toute relation
  4958	d'équivalence y soit effective, et tout morphisme s'y factorise en
  4959	épimorphisme effectif suivi par un monomorphisme effectif. Montrer
  4960	que dans $\Ind(C)$ les relations d'équivalence sont universelles si
  4961	et seulement si $C$ satisfait à la condition suivante:
  4962	\begin{itemize}
  4963	\item[%%
  4964	ST)] Pour tout objet $X$ de $C$ et tout sous-objet $Z$ de $X\times
  4965	X$ dans $C$, si on définit par récurrence la suite de sous-objets
  4966	$Z_n$ $(n\geq 0)$ de $X\times X$ dans $C$ par $Z_\circ=Z$,
  4967	$Z_n=\Sup(Z_n,s(Z_n),Z_n\circ Z_n)$ (le $\Sup$ pris dans l'ensemble
  4968	des sous-objets de $X\times X$, qui existe grâce aux hypothèses
  4969	faites sur C), alors la suite $(Z_n)_{n\geq 0}$ est
  4970	\emph{stationnaire.}
  4971	\end{itemize}
  4972	
  4973	\setcounter{enumi}{10}
  4974	\item Montrer que dans $\Ind((\Ens))$, les relations d'équivalence
  4975	ne sont pas nécessairement effectives.
  4976	
  4977	$^\ast$ NB. Pour des critères pour que $\Ind(C)$ soit un topos, \cf
  4978	VI 8.9.9.${}_\ast$
  4979	\end{enumerate}
  4980	\end{enonce*}
  4981	
  4982	\begin{enonce*}[remark]{Exercice 8.9.10}\etiquette[8.9.10]{I.exer8.9.10}
  4983	Soit\pageoriginale $C$ une $𝒰$-catégorie. Posons
  4984	$\Pro(C)=(\Ind(C^\circ))^\circ$ (\cf 8.11).
  4985	\begin{enumerate}
  4986	\renewcommand{\theenumi}{\alph{enumi}}
  4987	\renewcommand{\labelenumi}{\theenumi)}
  4988	\item Supposons que dans $C$ les sommes finies (\resp les petites
  4989	sommes) sont représentables. Montrer que si elles sont disjointes
  4990	(\cf II 4.5), alors $\Pro(C)$ satisfait la même %%
  4991	condition.
  4992	
  4993	\item Soit $J$ un petit ensemble, tel que les sommes indexées par
  4994	$J$ soient représentables dans $C$, donc aussi dans $\Pro(C)$. Soit
  4995	$(X(\alpha))_{\alpha\in J}$ une famille d'éléments de $\Pro(C)$,
  4996	avec $X(\alpha)=\textrm{« }\varprojlim_{I_\alpha}\textrm{ »\,}X_{i_\alpha}$.
  4997	Montrer que pour que la somme des $X(\alpha)$ dans $\Ind(C)$ soit
  4998	universelle, il suffit qu'il en soit de même pour chacune des
  4999	familles $\varprojlim_J X_{i_\alpha}$, où $(i_\alpha)_{a\in I}\in
  5000	\prod_{\alpha\in J}I_\alpha$. En conclure que si dans $C$ les sommes
  5001	de type $J$ sont universelles, pour qu'il en soit de même dans
  5002	$\Pro(C)$, il faut et il suffit que pour toute famille
  5003	$(X(\alpha))_{\alpha\in J}$ comme dessus, avec $I_\alpha=I$ pour
  5004	tout $\alpha\in J$, l'homomorphisme canonique dans $\Pro(C)$
  5005	$$
  5006	\textrm{« }\varprojlim_I\textrm{ »\,}\amalg_\alpha X(\alpha)_i\longleftarrow
  5007	\textrm{« }\varprojlim_{I^J}\textrm{ »\,} \amalg_{\alpha\in
  5008	J}X(\alpha)_{i_\alpha}
  5009	$$
  5010	soit un isomorphisme. En conclure que dans $\Pro((\Ens))$ les sommes
  5011	de type $J$ sont universelles si et seulement si $J$ est fini.
  5012	\end{enumerate}
  5013	\end{enonce*}
  5014	
  5015	\begin{enonce*}[remark]{Exercice 8.9.11}\etiquette[8.9.11]{I.exer8.9.11}
  5016	Soit $C$ une $𝒰$-catégorie.
  5017	\begin{enumerate}
  5018	\renewcommand{\theenumi}{\alph{enumi}}
  5019	\renewcommand{\labelenumi}{\theenumi)}
  5020	\item Soit $(X_i\rightarrow X)_{i\in I}$ une famille de morphismes
  5021	dans $C$. Montrer que pour qu'elle soit épimorphique dans $\Ind(C)$,
  5022	il suffit qu'il existe une sous-famille \emph{finie} qui soit
  5023	épimorphique dans $C$. Prouver que cette condition est également
  5024	nécessaire lorsqu'on suppose que dans $C$ toute famille finie de
  5025	morphismes de but $X_i$ se factorise en une famille
  5026	épimorphique,\pageoriginale suivie d'un monomorphisme effectif
  5027	(10.5). (Pour cette dernière assertion, soit $P$ l'ensemble des
  5028	      parties finies de $I$, ordonné par inclusion, et soit
  5029	      $X'=(X'_J)_{J\in P}$ le ind-objet formé des images des
  5030	      sous-familles finies de la famille donnée, enfin soit
  5031	      $X''=X\amalg_{X'}X=\textrm{« }\lim\textrm{ »\,}X\amalg_{X'_J}X$. Montrer
  5032	      que les deux morphismes canoniques $X\rightrightarrows X''$ sont
  5033	      distincts, mais coïncident sur les $X'_J$.)
  5034	
  5035	\item Supposons que dans $C$ toute famille finie de morphismes de
  5036	      même but se factorise en une famille épimorphique, suivie
  5037	      d'un monomorphisme effectif. Prouver que $\Ind(C)$ admet une
  5038	      petite sous-catégorie génératrice par épimorphismes
  5039	      (\hyperref[I.def7.1]{7.1}) si et seulement si $C$ admet une petite
  5040	      sous-catégorie
  5041	      $C'$ telle que tout objet de $C$ soit but d'une famille
  5042	      épimorphique \emph{finie} de source dans $C'$. Lorsqu'on
  5043	      suppose que $C$ admet une petite sous-catégorie
  5044	      génératrice (\hyperref[I.def7.1]{7.1}), alors $\Ind(C)$ admet une petite
  5045	      sous-catégorie génératrice par épimorphismes si et
  5046	      seulement si $C$ est équivalente à une petite
  5047	      catégorie. (Pour ce dernier énoncé, utiliser
  5048	      \hyperref[I.var7.5.2]{7.5.2}).
  5049	
  5050	\item $\Ind((\Ens))$ n'admet pas de petite sous-catégorie
  5051	  génératrice.
  5052	\end{enumerate}
  5053	\end{enonce*}
  5054	
  5055	\begin{enonce*}[remark]{Exercice 8.9.12}\etiquette[8.9.12]{I.8.9.12}
  5056	Soit $C$ une $𝒰$-catégorie, et considérons
  5057	$\Pro(C)=\Ind(C^\circ)^\circ$.
  5058	\begin{enumerate}
  5059	\renewcommand{\theenumi}{\alph{enumi}}
  5060	\renewcommand{\labelenumi}{\theenumi)}
  5061	\item Montrer que si $\Pro(C)$ admet une petite sous-catégorie $P'$
  5062	génératrice par épimorphismes (7.1), alors il existe une petite
  5063	sous-catégorie $C'$ de $C$ qui est génératrice par épimorphismes
  5064	dans $\Pro(C)$. (Prendre la catégorie des composants des objets de
  5065	$P'$.)
  5066	
  5067	\item Montrer\pageoriginale que la catégorie $\Pro((\Ens))$ n'a pas de petite
  5068	sous-catégorie génératrice par épimorphismes. (Si $C'$ est comme
  5069	dans a), choisir un ensemble $X$ dont le cardinal majore strictement
  5070	les cardinaux des éléments de $C'$, et considérer le pro-ensemble
  5071	$𝒳$ formé par les complémentaires dans $X$ des parties de
  5072	cardinal $<\card(X)$. Montrer que pour tout ensemble $Y$ non vide de
  5073	cardinal $<\card(X)$, on a $\Hom(Y,𝒳)=∅$, mais que
  5074	$𝒳$ n'est pas isomorphe au pro-ensemble constant
  5075	$∅$.)
  5076	\end{enumerate}
  5077	\end{enonce*}
  5078	
  5079	\subsection{Notions duales: pro-objets, foncteurs pro-représentables}\etiquette[8.10]{I.sec8.10}
  5080	
  5081	Soit $C$ une $𝒰$-catégorie. On appelle \emph{pro-objet} de
  5082	$C$ tout foncteur
  5083	$$
  5084	\displaylines{\rlap{$(8.10.1)$}\hfill \varphi:I^\circ\longrightarrow
  5085	C{\quoi},\hfill}
  5086	$$
  5087	où $I$ est une catégorie filtrante essentiellement petite (appelée
  5088	la catégorie d'indices), et comme d'habitude $I^\circ$ désigne la
  5089	catégorie opposée. Soit $𝒱$ un univers tel que
  5090	$𝒱\supset 𝒰$. On fera attention que les pro-objets
  5091	de $C$ indexés par des $I\in 𝒱$ sont en correspondance
  5092	biunivoque avec les ind-objets de $C^\circ$ indexés par des $I\in
  5093	𝒱$, en associant à tout tel ind-objet $\psi:I\rightarrow
  5094	C^\circ$ le pro-objet $\varphi=\psi^\circ:I^\circ\rightarrow C$.
  5095	S'inspirant de cette correspondance, on définit « par transport de
  5096	structure et renversement des flèches » la notion de morphisme
  5097	entre pro-objets de $C$ à partir de la notion analogue (8.2.4.2).
  5098	(8.2.5.1) pour les in-objets, et on définit en conséquence la
  5099	catégorie des pro-objets de $C$ indexés par des $I\in 𝒱$, et
  5100	un isomorphisme canonique:
  5101	$$
  5102	\displaylines{\rlap{$(8.10.2)$}\hfill
  5103	\Pro_{𝒱}(C,𝒰)\simeq \Ind_{𝒱}(C^\circ,
  5104	𝒰)^\circ{\quoi};\hfill}
  5105	$$
  5106	pour\pageoriginale $𝒱=𝒰$ on parle simplement de la
  5107	\emph{catégorie des pro-objets de} $C$, notée $\Pro_{𝒰}(C)$
  5108	ou simplement $\Pro(C)$, qui est donc définie par
  5109	$$
  5110	\displaylines{\rlap{$(8.10.3)$}\hfill
  5111	\Pro(C)=\Pro_{𝒰}(C,𝒰)\simeq
  5112	\Ind(C^\circ)^\circ{\quoi}.\hfill}
  5113	$$
  5114	Si deux pro-objets sont donnés sous \emph{forme indicielle}
  5115	\begin{equation}
  5116	𝒳=(X_i)_{i\in I}\quav 𝒴=(Y_j)_{j\in
  5117	J}{\quoi},\tag{8.10.4}
  5118	\end{equation}
  5119	la formule (8.2.5.1) prend ici la forme
  5120	\begin{equation}
  5121	\Hom(𝒳,𝒴)\simeq \varprojlim_j \varinjlim_i
  5122	\Hom(X_i,Y_j){\quoi}.\tag{8.10.5}
  5123	\end{equation}
  5124	Le foncteur (8.4.1) appliquée à $C^\circ$ donne, par passage aux
  5125	catégories opposées, un foncteur canonique (dénoté par
  5126	la même
  5127	lettre $c$ s'il n'y a pas risque de confusion), permettant
  5128	d'identifier $C$ à une sous-catégorie pleine de $\Pro(C)$:
  5129	\begin{equation}
  5130	c:C\longrightarrow\Pro(C){\quoi};\tag{8.10.6}
  5131	\end{equation}
  5132	c'est pour avoir un tel foncteur, et non $C\rightarrow
  5133	\Pro(C)^\circ$, qu'on a « renversé les flèches » dans la
  5134	définition des morphismes de pro-objets à partir de la définition
  5135	analogue pour les ind-objets. Les pro-objets dans l'image
  5136	essentielle de (8.10.6) s'appellent encore \emph{pro-objets
  5137	essentiellement constants.}
  5138	
  5139	Posons
  5140	\begin{equation}
  5141	\check{C}=(C^\circ)^{\widehat{\phantom{,}}}=\Hhom(C,(\Ens)){\quoi},\tag{8.10.7}
  5142	\end{equation}
  5143	alors le foncteur (8.2.4.7) peut être considéré
  5144	comme un foncteur
  5145	canonique pleinement fidèle
  5146	\begin{equation}
  5147	L:\Pro_{𝒱}(C,𝒰)\longrightarrow
  5148	\check{C}^\circ{\quoi},\tag{8.10.8}
  5149	\end{equation}\pageoriginale
  5150	et on trouve en particulier
  5151	\begin{equation}
  5152	L:\Pro(C)\hookrightarrow \check{C}^\circ{\quoi}.\tag{8.10.9}
  5153	\end{equation}
  5154	
  5155	Nous laisserons au lecteur le soin de traduire, au fur et mesure des
  5156	besoins, les résultats des sections précédentes et de celles qui
  5157	suivent du langage des ind-objets dans celui des pro-objets. Suivant
  5158	le contexte mathématique, c'est l'un ou l'autre langage qui est le
  5159	plus utile, sans compter les cas où les deux notions s'introduisent
  5160	simultanément, par exemple lorsqu'il y a lieu de considérer des
  5161	catégories complexes comme $\Pro(\Ind(C))$ ou $\Ind(\Pro(C))$. Nous
  5162	nous contentons de donner quelques indications supplémentaires, pour
  5163	fixer la terminologie et les notations, et préciser le « yoga ».
  5164	
  5165	
  5166	\setcounter{subsubsection}{9}
  5167	\subsubsection{}\etiquette[8.10.10]{I.sec8.10.10}
  5168	Un foncteur \emph{covariant $F:C\rightarrow (\Ens)$} est dit
  5169	\emph{foncteur pro-représentable} s'il est dans l'image essentielle
  5170	de (8.10.9) (ou 8.10.8, cela revient au même); la
  5171	sous-catégorie
  5172	pleine de $\check{C}$ formée de ces foncteurs est donc équivalente à
  5173	$\Pro(C)^\circ$. On préférera généralement travailler dans la
  5174	catégorie opposée, image essentielle \emph{en tant que
  5175	sous-catégorie de} (8.10.9), qui est donc équivalente à $\Pro(C)$
  5176	lui-même. En accord avec cet usage, il est souvent préférable de
  5177	regarder les foncteurs covariants
  5178	$$
  5179	F:C\longrightarrow (\Ens)
  5180	$$
  5181	comme des objets de $\Hhom(C,(\Ens))^\circ=\check{C}^\circ$, et
  5182	d'écrire en conséquence la catégorie des homomorphismes
  5183	$$
  5184	X\longrightarrow F\mbox{ dans } \check{C}^\circ{\quoi},
  5185	$$
  5186	\ie\pageoriginale des couples $(X,u)$ d'un objet $X$ de $C$ et d'un
  5187	$u\in F(X)$, comme
  5188	\begin{equation}
  5189	{}_{\displaystyle F}\backslash C{\quoi}  %%
  5190	\ .\tag{8.10.11}
  5191	\end{equation}
  5192	Avec cette convention, on trouvera donc un foncteur \emph{covariant}
  5193	(foncteur source)
  5194	\begin{equation}
  5195	{}_{\displaystyle F}\backslash C\rightarrow C{\quoi},\tag{8.10.12}
  5196	\end{equation}
  5197	et 3.4 se récrit sous la forme
  5198	\begin{equation}
  5199	F\xrightarrow{\sim}%%
  5200	\textrm{« }\varprojlim_{({}_{\displaystyle F}\backslash C)^\circ}\textrm{ »\,}
  5201	X\tag{8.10.13}
  5202	\end{equation}
  5203	
  5204	\setcounter{subsubsection}{13}
  5205	\subsubsection{}\etiquette[8.10.14]{I.sec8.10.14}
  5206	Le critère \hyperref[I.them8.3.3]{8.3.3} appliqué à $C^\circ$ devient
  5207	un critère de
  5208	proreprésentabilité pour $F:$ il faut et il suffit que
  5209	$({}_F\backslash C)^\circ$ soit filtrante et essentiellement petite,
  5210	cette dernière condition étant superflue si $C$ est équivalente à
  5211	une petite catégorie; si de plus dans $C$ les limites projectives
  5212	finies sont représentables, il revient au même de dire que $F$ y
  5213	commute \ie est \emph{exact à gauche.}
  5214	
  5215	\subsubsection{}\etiquette[8.10.15]{I.sec8.10.15}
  5216	Dans $\Pro(C)$ les petites limites projectives filtrantes sont
  5217	représentables et le foncteur (8.10.9) y commute; en d'autres
  5218	termes, ce foncteur transforme limites projectives en $\Pro(C)$ en
  5219	limites \emph{inductives} de $\check{C}=\Hhom(C,(\Ens))$, tout comme
  5220	son composé $C\rightarrow \check{C}^\circ$ avec (8.10.6), qui
  5221	partage avec lui la fâcheuse propriété d'être
  5222	contravariant quand on
  5223	la considère à valeurs dans $\check{C}$. On fera attention que les
  5224	foncteurs pro-représentables\pageoriginale de $C$ dans $(\Ens)$ sont
  5225	les foncteurs qui sont des petites limites \emph{inductives}
  5226	filtrantes de foncteurs représentables (et non des limites
  5227	projectives, comme la terminologie pourrait éventuellement le
  5228	suggérer).
  5229	
  5230	\subsection{Ind-adjoints et pro-adjoints}\etiquette[8.11]{I.sec8.11}
  5231	
  5232	\subsubsection{}\etiquette[8.11.1]{I.sec8.11.1}
  5233	Soit
  5234	\begin{equation}
  5235	f:C\longrightarrow C'\tag{8.11.1.1}
  5236	\end{equation}
  5237	un foncteur entre $𝒰$-catégories, d'où un foncteur
  5238	$F \sisi{→}{↦} F\circ f$
  5239	\begin{equation}
  5240	f^\ast:\widehat{C}'\longrightarrow\widehat{C}{\quoi}.\tag{8.11.1.2}
  5241	\end{equation}
  5242	On dit que $f$ \emph{admet un ind-adjoint} si le foncteur précédent
  5243	transforme foncteurs ind-représentables en foncteurs
  5244	ind-représentables. Comme $f^\ast$ commute aux limites inductives,
  5245	et que la sous-catégorie pleine de $C'$ formée des foncteurs
  5246	ind-représentables est stable par petites limites inductives
  5247	filtrantes, il revient au même de dire que $f$ admet un ind-adjoint,
  5248	ou que $f^\ast$ applique objets de $C$ (identifié à une
  5249	sous-catégorie pleine de $\widehat{C}$) dans des foncteurs
  5250	ind-représentables sur $C$, \ie
  5251	\begin{equation}
  5252	X\longmapsto \Hom(f(X),Y')\mbox{ est ind-représentable pour tout }
  5253	Y'\in \Ob C'{\quoi}.\tag{8.11.1.3}
  5254	\end{equation}
  5255	Une autre façon d'exprimer la condition que $f$ admette un
  5256	ind-adjoint est de dire qu'il existe un foncteur
  5257	\begin{equation}
  5258	g:\Ind(C')\longrightarrow \Ind(C)\tag{8.11.1.4}
  5259	\end{equation}
  5260	qui soit essentiellement induit par $f^\ast$, \ie tel qu'on ait un
  5261	isomorphisme de bifoncteurs
  5262	\begin{equation}
  5263	\Hom_{\Ind(C)}(X,g(𝒴'))\simeq
  5264	\Hom_{\Ind(C')}(f(X),𝒴'){\quoi},\tag{8.11.1.5}
  5265	\end{equation}
  5266	où\pageoriginale $X\in \Ob C$ et $Y'\in \Ob \Ind(C')$. En fait, il
  5267	suffit même d'avoir un foncteur
  5268	\begin{equation}
  5269	g_\circ:C'\longrightarrow \Ind(C)\tag{8.11.1.6}
  5270	\end{equation}
  5271	et un isomorphisme de bifoncteurs
  5272	\begin{equation}
  5273	\Hom_{\Ind(C)}(X,g_\circ(Y'))\simeq \Hom_{C'}(f(X),Y')\tag{8.11.1.7}
  5274	\end{equation}
  5275	en $X\in \Ob C$, $Y'\in \Ob C'$. Le foncteur $g$ (\resp $g_\circ$)
  5276	est évidemment déterminé à isomorphisme unique près par $f$, et
  5277	inversement on reconstitue $f$ à isomorphisme canonique près par la
  5278	connaissance de ce $g$ (\resp $g_\circ$). Il est clair sur
  5279	(8.11.1.5) que le foncteur $g$ commute aux limites inductives, et
  5280	qu'il « prolonge » le foncteur $g_\circ$; cela le détermine
  5281	donc à isomorphisme unique près en termes de $g_\circ$
  5282	(\hyperref[I.sec8.7.2]{8.7.2}). Le
  5283	foncteur $g$, et parfois aussi le foncteur $g_\circ$ qu'il prolonge,
  5284	est appelé le \emph{foncteur ind-adjoint de} $f$ (inutile ici de
  5285	préciser: à droite, car l'autre, s'il existe, s'appellera le
  5286	pro-adjoint, \cf \hyperref[I.sec8.11.5]{8.11.5} plus bas).
  5287	
  5288	Bien entendu, lorsque $f$ admet un adjoint à droite
  5289	$$
  5290	f^{\ad}:C'\longrightarrow C{\quoi},
  5291	$$
  5292	il admet un ind-adjoint $g$, et celui-ci est isomorphe canoniquement
  5293	au prolongement canonique de $f$ aux ind-objets:
  5294	\begin{equation}
  5295	g\simeq\Ind(f^{\ad}){\quoi}.\tag{8.11.1.8}
  5296	\end{equation}
  5297	La notion de ind-adjoint est donc une généralisation naturelle de la
  5298	notion d'adjoint à droite.
  5299	
  5300	Considérons\pageoriginale maintenant le prolongement canonique
  5301	\begin{equation}
  5302	\Ind(f):\Ind(C)\longrightarrow\Ind(C'){\quoi}.\tag{8.11.1.9}
  5303	\end{equation}
  5304	On a alors:
  5305	
  5306	\begin{enonce*}{Proposition 8.11.2}\etiquette[8.11.2]{I.prop8.11.2}
  5307	Pour que le foncteur $f:C\rightarrow C'$ admette un ind-adjoint, il
  5308	faut et il suffit que le foncteur $\Ind(f)$ $(8.11.1.9)$ admette un
  5309	adjoint à droite. Ce dernier est canoniquement isomorphe au
  5310	ind-adjoint $(8.11.1.4)$.
  5311	\end{enonce*}
  5312	
  5313	Le suffisance et la dernière assertion sont triviales sur la formule
  5314	de ind-adjonction (8.11.1.5). Pour la nécessité, on
  5315	note que par
  5316	passage à la limite projective sur cette formule sur des $X_i$, on
  5317	déduit un isomorphisme d'adjonction en les pro-objets
  5318	$𝒳=(X_i)_{i\in I}$ et $𝒴'$:
  5319	\begin{equation}
  5320	\Hom_{\Ind(C)}(𝒳, g(𝒴'))\simeq
  5321	\Hom_{\Ind(C')}(\Ind(f)(𝒳),𝒴'){\quoi},\tag{8.11.2.1}
  5322	\end{equation}
  5323	\cqfd
  5324	
  5325	\begin{enonce*}{Corollaire 8.11.3}\etiquette[8.11.3]{I.cor8.11.3}
  5326	Si $f$ admet un ind-adjoint, alors $\Ind(f)$ commute aux limites
  5327	inductives, et le ind-adjoint $g$ commute aux limites projectives.
  5328	\end{enonce*}
  5329	
  5330	\begin{enonce*}{Proposition 8.11.4}\etiquette[8.11.4]{I.prop8.11.4}
  5331	Pour que le foncteur $f:C\rightarrow C'$ admette un ind-adjoint, il
  5332	faut que $f$ soit exact à \sisi{gauche}{droite}, et cette condition est suffisante
  5333	lorsque $C$ est équivalente à une petite catégorie.
  5334	\end{enonce*}
  5335	
  5336	Cela résulte du critère (8.11.1.3), et de \hyperref[I.prop8.3.1]{8.3.1}
  5337	et \hyperref[I.them8.3.3]{8.3.3} (iv).
  5338	
  5339	Pour un autre critère en termes de la notion de foncteur accessible,
  5340	\cf 8.13.3.
  5341	
  5342	\setcounter{subsubsection}{4}\etiquette[8.11.5]{I.sec8.11.5}
  5343	\subsubsection{}
  5344	Considérons\pageoriginale maintenant le foncteur $F\mapsto F\circ f$
  5345	\begin{equation}
  5346	f^{\circ\ast}:\check{C}'\longrightarrow \check{C}\tag{8.11.5.1}
  5347	\end{equation}
  5348	induit par \sisi{$F$}{$f$}. On dira que $f$ \emph{admet un pro-adjoint} si le
  5349	foncteur précédent applique foncteur pro-représentable en foncteur
  5350	pro-représentable, \ie s'il existe un foncteur (appelé
  5351	\emph{foncteur pro-adjoint de $f$})
  5352	\begin{equation}
  5353	g:\Pro(\sisi{C}{C'})\longrightarrow \Pro(\sisi{C'}{C}){\quoi},\tag{8.11.5.2}
  5354	\end{equation}
  5355	et un isomorphisme de bifoncteurs
  5356	\begin{equation}
  5357	\Hom_{\Pro(C)}(g(𝒴'),𝒳)\simeq
  5358	\Hom_{\Pro(C')}(𝒴'\sisi{)}{},f(𝒳)){\quoi}.\tag{8.11.5.3}
  5359	\end{equation}
  5360	Bien entendu, dire que $f$ admet un pro-adjoint signifie que
  5361	$f^\circ$ admet un ind-adjoint, de sorte que les notions et
  5362	résultats pour les ind-adjoints se traduisent trivialement en termes
  5363	de pro-adjoints. Signalons seulement que $f$ admet un pro-adjoint si
  5364	et seulement si $\Pro(f):\Pro(C)\rightarrow \Pro(C')$ admet un
  5365	adjoint \emph{à \sisi{droite}{gauche}}, et que dans ce cas le foncteur $g$
  5366	précédent est un tel adjoint à \sisi{droite}{gauche} de
  5367	$\Pro(\sisi{C}{f})$; et qu'il faut
  5368	pour ceci que $f$ soit exact à \emph{gauche}, cette condition étant
  5369	également suffisante lorsque $C$ est équivalente à une petite
  5370	catégorie. Dans ce cas, $f$ est donc exact si et seulement si il
  5371	admet à la fois un ind-adjoint et un pro-adjoint.
  5372	
  5373	\begin{enonce*}[remark]{Exemple 8.11.6}\etiquette[8.11.6]{I.exem8.11.6}
  5374	Considérons le cas d'un foncteur
  5375	$$
  5376	f:C\longrightarrow (\Ens){\quoi}.
  5377	$$
  5378	Si ce foncteur admet un pro-adjoint, il est pro-représentable, et la
  5379	réciproque est vraie si et seulement si la sous-catégorie pleine de
  5380	$\check{C}$ formée des foncteurs pro-représentables est stable par
  5381	petits produits; c'est\pageoriginale le cas en particulier si $C$
  5382	est équivalente à une petite catégorie (\hyperref[I.sec8.10.14]{8.10.14}) ou si
  5383	dans $C$ les petits produits sont représentables (\hyperref[I.prop8.9.5]{8.9.5}
  5384	b) appliquée à $C^\circ$). À peu de choses près, on peut donc dire
  5385	que pour un foncteur $f:C\rightarrow C'$, la notion d'existence d'un
  5386	pro-adjoint est la \emph{généralisation naturelle de la notion de
  5387	pro-représentabilité de $f$}, qui est définie lorsque $C'=(\Ens)$.
  5388	\end{enonce*}
  5389	
  5390	\subsection{Ind-objets et pro-objets stricts. Application à un
  5391	critère de représentabilité}\etiquette[8.12]{I.sec8.12}
  5392	
  5393	\subsubsection{}\etiquette[8.12.1]{I.sec8.12.1}
  5394	Soit
  5395	$$
  5396	𝒳=(X_i)_{i\in I}
  5397	$$
  5398	un ind-objet de la $𝒰$-catégorie $C$, et soit $F$ le
  5399	préfaisceau qu'il ind-représente. On voit alors aussitôt qu'il
  5400	revient au même de dire que les morphismes canoniques
  5401	$$
  5402	X_i\longrightarrow F=\textrm{« }\varinjlim_i\textrm{ »\,}X_i
  5403	$$
  5404	sont des monomorphismes de $\Ind(C)$ (ou, ce qui revient au même, de
  5405	$\widehat{C}$, \ie des monomorphismes de foncteurs « argument par
  5406	argument »), ou de dire que pour toute flèche $i\rightarrow j$ de
  5407	$I$, la flèche de transition correspondante
  5408	$$
  5409	X_i\longrightarrow X_j
  5410	$$
  5411	est un monomorphisme. Lorsque ces conditions sont remplies, et si de
  5412	plus $I$ est une catégorie ordonnée, on dit que $𝒳$ est un
  5413	\emph{ind-objet strict}. On\pageoriginale notera que cette condition
  5414	n'est pas invariante par isomorphisme de ind-objets; un ind-objet
  5415	sera appelé \emph{essentiellement strict} s'il est isomorphe à un
  5416	ind-objet strict. Un préfaisceau sera appelé \emph{strictement
  5417	ind-représentable} s'il est ind-représentable par un ind-objet
  5418	strict; donc $F=\textrm{« }\varinjlim_i\textrm{ »\,}X_i$ est strictement
  5419	ind-représentable si et seulement si $𝒳=(X_i)_{i\in I}$ est
  5420	essentiellement strict.
  5421	
  5422	\paragraph{}\etiquette[8.12.1.1]{I.8.12.1.1}
  5423	Soit $F$ un préfaisceau sur $C$, et considérons la sous-catégorie
  5424	pleine de $C_{/F}$ formée des flèches $X\rightarrow F$ de source
  5425	dans $C$ qui sont des monomorphismes. On l'appellera la
  5426	\emph{catégorie des sous-foncteurs représentables de $F$}; c'est la
  5427	catégorie associée à l'ensemble ordonné des sous-foncteurs
  5428	représentables de $F$, ordonné par l'ordre induit de celui de
  5429	l'ensemble des sous-objets de $F$. Il résulte alors aussitôt des
  5430	définitions:
  5431	
  5432	\begin{enonce*}{Proposition 8.12.2}\etiquette[8.12.2]{I.prop8.12.2}
  5433	Pour que le préfaisceau $F$ sur $C$ soit strictement
  5434	ind-représentable, il faut et il suffit que la catégorie (ordonnée)
  5435	$I$ des sous-foncteurs représentables de $F$ soit filtrante et
  5436	essentiellement petite et que l'on ait
  5437	\begin{equation*}
  5438	\varinjlim_I X_i \xrightarrow{\sim}F{\quoi}.\tag{8.12.2.1}
  5439	\end{equation*}
  5440	Lorsque pour tout objet $X$ de $C$, l'ensemble des sous-objets de
  5441	$X$ dans $C$ est petit (par exemple si $C$ admet une petite
  5442	sous-catégorie génératrice $(7.4)$), cela implique que la catégorie
  5443	des sous-foncteurs représentables est même petite.
  5444	\end{enonce*}
  5445	
  5446	\setcounter{subsubsection}{2} \setcounter{paragraph}{1}
  5447	\paragraph{}\etiquette[8.12.2.2]{I.sec8.12.2.2}
  5448	Si\pageoriginale $F$ est strictement ind-représentable, il y a donc
  5449	une façon privilégiée de le ind-représenter par un ind-objet, et ce
  5450	dernier est même un ind-objet strict: on prend la représentation
  5451	(8.12.2.1). Un ind-objet strict $𝒴$ est dit
  5452	\emph{saturé} s'il est isomorphe à un ind-objet de la forme $(X_i)$
  5453	figurant dans (8.12.2.1), pour $F$ convenable; donc il
  5454	existe à isomorphisme unique près, un seul ind-objet strict saturé
  5455	isomorphe au ind-objet strict donné, savoir celui envisagé dans
  5456	\hyperref[I.prop8.12.2]{8.12.2}, en prenant $F=\varinjlim 𝒴$ (limite dans
  5457	$\widehat{C}$).
  5458	
  5459	\paragraph{}\etiquette[8.12.2.3]{I.sec8.12.2.3}
  5460	Supposons qu'on sache déjà que l'on puisse trouver une petite partie
  5461	cofinale dans l'ensemble $\Ob C_{/F}$, ce qui est le cas en
  5462	particulier si $F$ est ind-représentable; alors il s'ensuit que la
  5463	même condition est vérifiée dans la sous-catégorie pleine $I$
  5464	envisagée dans \hyperref[I.prop8.12.2]{8.12.2}, donc on peut dans le critère
  5465	\hyperref[I.prop8.12.2]{8.12.2} omettre la condition que $I$ soit essentiellement
  5466	petite.
  5467	
  5468	\subsubsection{}\etiquette[8.12.3]{I.sec8.12.3}
  5469	Soit $F$ un préfaisceau sur $C$, et considérons un objet
  5470	$$
  5471	u:X\longrightarrow F
  5472	$$
  5473	de $C_{/F}$. On dit que $u$ (ou le couple $(X,u)$) est
  5474	\emph{minimal} si pour toute factorisation de $u$ en
  5475	\begin{equation}
  5476	X\xrightarrow{p}X'\xrightarrow{u'}F{\quoi},\tag{8.12.3.1}
  5477	\end{equation}
  5478	avec $p$ un épimorphisme strict (\hyperref[I.sec10.2]{10.2}), $p$ est un isomorphisme.
  5479	Considérant $u$ comme un objet de $F(X)$ (1.4), dire que $u$ est
  5480	minimal signifie donc que\pageoriginale tout épimorphisme strict
  5481	$p:X\rightarrow X'$ tel que $u\in \Iim(F(p):F(X')\rightarrow F(X))$
  5482	est un isomorphisme. Cette notion s'éclaire par la partie a) du
  5483	lemme suivant:
  5484	
  5485	\begin{enonce*}{Lemme 8.12.4}\etiquette[8.12.4]{I.lem8.12.4}
  5486	Soit $F$ un préfaisceau sur $C$.
  5487	\begin{enumerate}
  5488	\renewcommand{\theenumi}{\alph{enumi}}
  5489	\renewcommand{\labelenumi}{\theenumi)}
  5490	\item Supposons que $F$ transforme conoyaux en noyaux et considérons
  5491	un morphisme
  5492	$$
  5493	u:X\longrightarrow F{\quoi},
  5494	$$
  5495	avec $X\in \Ob C$. Pour que $u$ soit un monomorphisme, il suffit,
  5496	lorsque dans $C$ les conoyaux de doubles flèches sont
  5497	représentables, que $u$ soit minimal; cette condition est également
  5498	nécessaire si dans $C$ les produits fibrés sont représentables.
  5499	
  5500	\item Supposons que dans $C$ les $\varinjlim$ finies soient
  5501	représentables. Pour que %%
  5502	la sous-catégorie des sous-foncteurs représentables de $F$
  5503	(\hyperref[I.8.12.1.1]{8.12.1.1}) soit filtrante (pas nécessairement petite) et ait
  5504	pour limite inductive $F$, il suffit que $F$ soit exact à gauche et
  5505	que tout morphisme $u:X\rightarrow F$, avec $X\in \Ob C$, se
  5506	factorise en
  5507	$$
  5508	X\rightarrow X'\xrightarrow{u'}F{\quoi},
  5509	$$
  5510	avec $u'$ minimal; cette condition est également nécessaire si dans
  5511	$C$ les produits fibrés sont représentables.
  5512	
  5513	\item Supposons que $C$ admette une petite sous-catégorie
  5514	génératrice (\hyperref[I.def7.1]{7.1}), et que $I$ la catégorie des sous-foncteurs
  5515	représentables de $F$ soit filtrante, alors $I$ est petite.
  5516	\end{enumerate}
  5517	\end{enonce*}
  5518	
  5519	\noindent \emph{Démonstration}.\pageoriginale
  5520	\begin{enumerate}
  5521	\renewcommand{\theenumi}{\alph{enumi}}
  5522	\renewcommand{\labelenumi}{\theenumi)}
  5523	\item Supposons $u$ minimal, prouvons que $u$ est un monomorphisme,
  5524	\ie que pour toute double-flèche $v$, $v':Y\rightrightarrows X$
  5525	telle que $uv=uv'$, on a $v=v'$. En effet, si $X'=\Coker(v,v')$,
  5526	alors $u$ se factorise en $X\xrightarrow{p}X'\xrightarrow{u'}F$ ($F$
  5527	transformant conoyaux en noyaux), et comme $p$ est un épimorphisme
  5528	strict par construction, il s'ensuit que $p$ est un isomorphisme \ie
  5529	$v=v'$. Supposons que $u$ est un monomorphisme, prouvons qu'il est
  5530	minimal. Considérons une factorisation (8.12.3.1); comme $p$ est un
  5531	épimorphisme strict, c'est le conoyaux de la double flèche canonique
  5532	$v$, $v':\fprod{X}{X}{X'}\rightrightarrows X$, et comme
  5533	$(u'p)v=(u'p)v'$ et que $u'p=u$ est un monomorphisme, on a $v=v'$
  5534	donc $p$ est un isomorphisme.
  5535	
  5536	\item Suffisance: Comme $F$ est exact à gauche, $C_{/F}$ est
  5537	filtrante. Comme on sait que $F=\varinjlim_{C_{/F}}X$, on est ramené
  5538	par \hyperref[I.prop8.1.3]{8.1.3} c) à prouver que la sous-catégorie pleine $I$
  5539	de $C_{/F}$ des sous-objets de $F$ est cofinale dans $C_{/F}$; or en
  5540	vertu du « il suffit » dans a), c'est ce qu'assure l'hypothèse
  5541	que tout objet de $C_{/F}$ est majoré par un objet « minimal ».
  5542	Nécessité: Comme toute limite inductive filtrante de foncteurs
  5543	exacts à droite est itou, la première condition est trivialement
  5544	nécessaire. La deuxième résulte alors du « il faut » dans a).
  5545	
  5546	\item Un sous-foncteur représentable $X\hookrightarrow F$ de $F$ est
  5547	connu quand on connaît la sous-catégorie pleine $C'_{/X}$ de
  5548	$C'_{/F}$, où $C$ est une petite sous-catégorie génératrice fixée de
  5549	$C$. (Utiliser l'hypothèse $I$ filtrante.) Comme $C'_{/F}$ est
  5550	petite, l'ensemble de ses sous-catégories pleine est petit, d'où le
  5551	conclusion.
  5552	\end{enumerate}
  5553	
  5554	\begin{enonce*}{Proposition 8.12.5}\etiquette[8.12.5]{I.prop8.12.5}
  5555	Soit\pageoriginale $C$ une $𝒰$-catégorie où les limites
  5556	inductives finies sont représentables, et admettant une petite
  5557	sous-catégorie génératrice (\hyperref[I.def7.1]{7.1}). Soit $F$ un
  5558	préfaisceau sur $C$.
  5559	Pour que $F$ soit strictement ind-représentable, il suffit qu'il
  5560	satisfasse les deux conditions suivantes, et celles-ci sont
  5561	également nécessaires si dans $C$ les produits fibrés sont
  5562	représentables:
  5563	\begin{enumerate}
  5564	\renewcommand{\theenumi}{\alph{enumi}}
  5565	\renewcommand{\labelenumi}{\theenumi)}
  5566	\item $F$ est exact à gauche.
  5567	
  5568	\item Tout couple $(X,u)$, avec $X\in \Ob C$ et $u\in F(X)$, est
  5569	majoré dans $C_{/F}$ par un couple minimal $(8.12.3)$, \ie il existe
  5570	un couple minimal $(X',u')$ $(X'\in \Ob C, u'\in F(X))$ et un
  5571	morphisme $f:X\rightarrow X'$ tel que $u=F(f)(u')$.
  5572	\end{enumerate}
  5573	\end{enonce*}
  5574	
  5575	La suffisance résulte de \hyperref[I.prop8.12.2]{8.12.2} et de \hyperref[I.lem8.12.4]{8.12.4}
  5576	b), c), la nécessité de \hyperref[I.prop8.3.1]{8.3.1} et de \hyperref[I.lem8.12.4]{8.12.4}
  5577	a).
  5578	
  5579	\begin{enonce*}[remark]{Remarque 8.12.6}\etiquette[8.12.6]{I.rem8.12.6}
  5580	Lorsque $F$ transforme sommes amalgamées de $C$ en produits fibrés,
  5581	on voit aussitôt que pour un couple $(X,u)$ donné comme dans b),
  5582	l'ensemble des quotients stricts $X'$ de $X$ tels que $u\in
  5583	\Iim(F(X')\rightarrow F(X))$ est filtrant décroissant, ce qui
  5584	implique que si l'ensemble des quotients stricts de $X$ est
  5585	artinien, alors la condition b) de \hyperref[I.prop8.12.5]{8.12.5} est
  5586	automatiquement satisfaite. Si donc la condition précédente sur $X$
  5587	est satisfaite pour tout objet $X$ de $C$ (on dit aussi alors,
  5588	parfois, que les objets de $C^\circ$ sont \emph{artiniens}), alors
  5589	il résulte de 8.12.5 que $F$ est strictement ind-représentable si et
  5590	seulement si $F$ est exact à gauche; dans ce cas, $F$ est donc
  5591	strictement ind-représentable dès qu'il est ind-représentable
  5592	(\hyperref[I.prop8.3.1]{8.3.1}).
  5593	\end{enonce*}
  5594	
  5595	\begin{enonce*}{Corollaire 8.12.7}\etiquette[8.12.7]{I.cor8.12.7}
  5596	Soit\pageoriginale $C$ une $𝒰$-catégorie où les petites
  5597	limites projectives sont représentables, et admettant une petite
  5598	sous-catégorie cogénératrice (\hyperref[I.sec7.9]{7.9},
  5599	\hyperref[I.exem7.13]{7.13}). Alors un foncteur
  5600	$F:C\rightarrow (\Ens)$ est représentable si et seulement si $F$
  5601	commute aux petites limites projectives.
  5602	\end{enonce*}
  5603	La nécessité est claire. Pour la suffisance, on applique
  5604	\hyperref[I.prop8.12.5]{8.12.5}  à la catégorie opposée $C^\circ$. Pour prouver
  5605	d'abord que $F$ est (strictement) pro-représentable, on est ramené à
  5606	prouver que tout couple $(X,u)$, $X\in \Ob C$ et $u\in F(X)$, est
  5607	majoré par un couple « minimal ». Or la famille $(X_i)_{i\in I}$
  5608	des sous-objets \emph{stricts} de $X$ tels que $u\in
  5609	\Iim(F(X_i)\rightarrow F(X))$ est petite (\hyperref[I.cor7.5]{7.5} sous forme
  5610	duale). Grâce au fait que $F$ est exact à gauche, elle est
  5611	co-filtrante (\hyperref[I.rem8.12.6]{8.12.6}), et grâce au fait que $F$ commute
  5612	aux petites limites projectives on voit que $X'=\varprojlim_i X_i$
  5613	est un plus petit objet de cette famille. (Utiliser le fait que, $F$
  5614	étant exact à gauche, transforme monomorphismes en monomorphismes.)
  5615	Si $u'\in F(X')$ est l'unique élément dont l'image dans $F(X)$ est
  5616	$u$, on voit alors que $(X',u')$ est un couple minimal majorant
  5617	$(X,u)$. Cela prouve que $F$ est pro-représentable, et la conclusion
  5618	résulte alors du
  5619	
  5620	\begin{enonce*}{Lemme 8.12.8.1}\etiquette[8.12.8.1]{I.lem8.12.8.1}
  5621	Soit $F:C\rightarrow(\Ens)$ un foncteur, où $C$ est une
  5622	$𝒰$-catégorie où les petites limites projectives sont
  5623	représentables. Pour que $F$ soit représentable, il faut et il
  5624	suffit qu'il soit pro-représentable et qu'il commute aux petites
  5625	limites projectives.
  5626	\end{enonce*}
  5627	
  5628	La\pageoriginale nécessité est claire, prouvons la suffisance. Dire
  5629	que $F$ est représentable signifie évidemment que ${}_F\backslash C$
  5630	admet un élément initial (en fait, les objets initiaux de
  5631	${}_F\backslash C$ sont précisément les isomorphismes $F\rightarrow
  5632	X$, \ie les données de représentation pour $F$). Or $F$ étant
  5633	proreprésentable, $({}_F\backslash C^\circ)$ est filtrante et
  5634	équivalente à une petite catégorie, donc
  5635	$X=\varprojlim_{({}_F\backslash C)^\circ} X$ est représentable dans
  5636	$C$, et comme $F$ commute à la limite envisagée, il s'ensuit que $X$
  5637	est un objet initial de ${}_F\backslash C$, \cqfd
  5638	
  5639	\begin{enonce*}{Corollaire 8.12.8}\etiquette[8.12.8]{I.cor8.12.8}
  5640	Les hypothèses sur $C$ étant celles de $8.12.7$, soit
  5641	$f:C\rightarrow C'$ un foncteur de $C$ dans une
  5642	$𝒰$-catégorie $C'$. Pour que $f$ admette un adjoint à
  5643	gauche, il faut et il suffit que $f$ commute aux petites limites
  5644	projectives.
  5645	\end{enonce*}
  5646	
  5647	Cela se ramène en effet trivialement à \hyperref[I.cor8.12.7]{8.12.7}, en appliquant cet
  5648	énoncé aux foncteurs composés de la forme $X\mapsto \Hom(%%
  5649	Y',f(X))$.
  5650	
  5651	\begin{enonce*}[remark]{Exemples 8.12.9}\etiquette[8.12.9]{I.exer8.12.9}
  5652	Comme on a signalé dans \hyperref[I.exem7.13]{7.13}, les hypothèses sur
  5653	$E$ de \hyperref[I.cor8.12.7]{8.12.7} et
  5654	\hyperref[I.cor8.12.8]{8.12.8} sont vérifiées si $C$ est la catégorie des
  5655	$𝒰$-faisceaux d'ensembles sur un espace topologique $X\in
  5656	𝒰$ (ou plus généralement, sur un $𝒰$-site (II
  5657	(3.0.2)). Donnons un exemple instructif \footnote[1]{(dû à H.
  5658	BASS).} qui montre que l'hypothèse d'existence d'une petite
  5659	sous-catégorie cogénératrice $D$ de $C$ n'est pas surabondante dans
  5660	\hyperref[I.cor8.12.7]{8.12.7}. Prenons pour $C$ la catégorie des groupes
  5661	éléments de
  5662	$𝒰$. Soit $J$ l'ensemble des classes d'isomorphie de groupes
  5663	simples $\in C$, choisissons pour tout $j\in J$ un groupe simple
  5664	$G_j$ dans la classe de $j$, et soit $I$ l'ensemble ordonné filtrant
  5665	des parties $𝒰$-petites de $J$, et pour $i\in I$, soit
  5666	$X_i=\amalg_{j\in i}G_j$. Les $X_i$ forment alors un\pageoriginale
  5667	système projectif $(X_i)_{i\in I}$ dans $E$, à morphismes de
  5668	transition des épimorphismes, et le foncteur correspondant
  5669	\begin{equation}
  5670	F(X)=\varinjlim_i \Hom(X_i,X)\tag{$\ast$}
  5671	\end{equation}
  5672	prend ses valeurs dans $(𝒰\text{-}\Ens)$ (bien que
  5673	l'ensemble d'indices $I$ n'ait évidemment pas un cardinal $\in
  5674	𝒰$); plus précisément, montrons que pour toute petite
  5675	sous-catégorie pleine $C_\circ$ de $C$, il existe un $i_\circ\in I$
  5676	tel que la restriction de $F$ à $C_\circ$ soit représentable par
  5677	$X_{i_\circ}$, ce qui prouvera à la fois que $F$ est à valeurs dans
  5678	$𝒰\text{-}\Ens$, et qu'il commute aux petites limites
  5679	projectives. Pour prouver notre assertion, il suffit de noter que
  5680	pour tout $X\in \Ob C_\circ$, le cardinal de l'ensemble $J(X)$ des
  5681	$j\in J$ tels qu'il existe un homomorphisme non trivial de $G_j$
  5682	dans $X$ est nécessairement petit, puisque un tel morphisme est
  5683	nécessairement un monomorphisme ($G_j$ étant simple); par suite, si
  5684	$i_\circ$ est la partie de $J$ réunion des $J(X)$ pour $X\in \ob
  5685	C_\circ$, $i_\circ$ est petit \ie $i_\circ\in I$, et il fait
  5686	l'affaire. D'autre part il est clair que $F$ n'est pas
  5687	représentable, puisque on a $\card I\not\in 𝒰$. De ceci et
  5688	de 8.12.7 on conclut donc que \emph{la catégorie $C$ des groupes
  5689	$\in 𝒰$ n'admet pas une petite sous-catégorie pleine
  5690	cogénératrice.} Comme l'objet $\ZZ$ de $C$ est d'autre part un
  5691	générateur, il résulte alors de la démonstration de
  5692	\hyperref[I.cor7.12]{7.12} qu'il
  5693	existe un groupe $G$ à deux générateurs qui ne se plonge
  5694	pas dans un
  5695	objet injectif de la catégorie $C$ des groupes $\in 𝒰$. Il
  5696	semble d'ailleurs plausible que $C$ n'admette pas d'autre objet
  5697	injectif que les groupes unités.
  5698	\end{enonce*}
  5699	
  5700	\subsection{Foncteurs proreprésentables et foncteurs
  5701	  accessibles}\etiquette[8.13]{I.sec8.13}\pageoriginale
  5702	
  5703	\subsubsection{}\etiquette[8.13.1]{I.sec8.13.1}
  5704	Dans le présent numéro, nous utilisons quelques notions et résultats
  5705	du paragraphe suivant, et notamment \hyperref[I.prop9.11]{9.11} et
  5706	\hyperref[I.prop9.13]{9.13}, pour obtenir un critère de proreprésentabilité que
  5707	nous utiliserons (incidemment) dans IV 9.16. $C$ désigne par la
  5708	suite une $𝒰$-catégorie satisfaisant la condition $L$ de
  5709	\hyperref[I.def9.1]{9.1} b), cette condition étant remplie par exemple si dans
  5710	$C$ les petites limites inductives filtrantes sont représentables.
  5711	
  5712	\begin{enonce*}{Proposition 8.13.2}\etiquette[8.13.2]{I.prop8.13.2}
  5713	Soient $C$ comme ci-dessus, et
  5714	$$
  5715	f:C\longrightarrow(\Ens)
  5716	$$
  5717	un foncteur.
  5718	\begin{enumerate}
  5719	\renewcommand{\theenumi}{\alph{enumi}}
  5720	\renewcommand{\labelenumi}{\theenumi)}
  5721	\item Supposons que chaque objet de $C$ est accessible (\hyperref[I.def9.3]{9.3}). Si $f$
  5722	est pro-représentable, $f$ est accessible (\hyperref[I.def9.2]{9.2}) et
  5723	exact à gauche.
  5724	
  5725	\item Supposons que dans $C$ les limites projectives finies soient
  5726	représentables, et que $C$ admette une filtration cardinale
  5727	(\hyperref[I.def9.12]{9.12}). Si $f$ est accessible et exact à gauche, alors $f$ est
  5728	proreprésentable.
  5729	\end{enumerate}
  5730	\end{enonce*}
  5731	
  5732	\noindent \emph{Démonstration.}
  5733	\begin{enumerate}
  5734	\renewcommand{\theenumi}{\alph{enumi}}
  5735	\renewcommand{\labelenumi}{\theenumi)}
  5736	\item L'hypothèse sur $C$ signifie que les foncteurs covariants
  5737	représentables de $C$ dans $(\Ens)$ sont accessibles. Il en est donc
  5738	de même de toute petite limite inductive de tels foncteurs
  5739	(\hyperref[I.prop9.6]{9.6} (i)), donc aussi de tout foncteur
  5740	pro-représentable.
  5741	
  5742	\item En vertu de \hyperref[I.them8.3.3]{8.3.3} (iii), il reste à prouver que dans
  5743	$\Ob({}_F\backslash C)^\circ$ il y a une petite sous-catégorie
  5744	cofinale. Or par hypothèse il existe un cardinal $\Pi$ tel que $f$
  5745	soit $\Pi$-accessible. Soit alors $(X,u)$,\pageoriginale $u\in
  5746	F(X)$, un objet de ${}_F\backslash C$. Avec les notations de 9.12 c)
  5747	on a alors $X=\varinjlim_i X_i$, avec $I$ filtrant grand devant
  5748	$\Pi$ et les $X_i$ dans $C'=\Filt^{\Pi}(C)$, d'où
  5749	$F(X)\xleftarrow{\sim}\varinjlim_i F(X_i)$. Cela montre que la
  5750	petite sous-catégorie $({}_F\backslash C')^\circ$ est cofinale dans
  5751	$({}_F\backslash C)^\circ$, et achève la démonstration.
  5752	\end{enumerate}
  5753	
  5754	\begin{enonce*}{Corollaire 8.13.3}\etiquette[8.13.3]{I.cor8.13.3}
  5755	Soit $C$ une $𝒰$-catégorie satisfaisant aux conditions
  5756	suivantes:
  5757	\begin{enumerate}
  5758	\renewcommand{\theenumi}{\alph{enumi}}
  5759	\renewcommand{\labelenumi}{\theenumi)}
  5760	\item Dans $C$ les limites projectives finies sont représentables.
  5761	
  5762	\item Dans $C$ les petites limites inductives filtrantes sont
  5763	représentables.
  5764	
  5765	\item Le foncteur $\Ker$ sur la catégorie des doubles flèches de
  5766	$C$ est accessible (par exemple, il commute aux petites $\varinjlim$
  5767	filtrantes).
  5768	
  5769	\item Tout épimorphisme strict de $C$ est strict universel $(10.2)$.
  5770	
  5771	\item Il existe une petite sous-catégorie de $C$ génératrice par
  5772	épimorphismes stricts $(7.1)$.
  5773	\end{enumerate}
  5774	Sous ces conditions, un foncteur $f:C\rightarrow (\Ens)$ est
  5775	proreprésentables si et seulement si il est exact à gauche et
  5776	accessible (\hyperref[I.def9.2]{9.2}).
  5777	\end{enonce*}
  5778	
  5779	En effet, les conditions de \hyperref[I.prop8.13.2]{8.13.2} a) et b) sur $C$ sont
  5780	vérifiées, en vertu de \hyperref[I.prop9.11]{9.11} et \hyperref[I.prop9.13]{9.13}
  5781	respectivement.
  5782	
  5783	
  5784	\begin{enonce*}{Corollaire 8.13.4}\etiquette[8.13.4]{I.cor8.13.4}
  5785	Soit $f:C\rightarrow C'$ un foncteur entre $𝒰$-catégories
  5786	satisfaisant aux conditions {\rm a)} à {\rm e)} de \hyperref[I.cor8.13.3]{8.13.3}. Pour
  5787	que $f$ admette un pro-adjoint, il faut et il suffit que $f$ soit
  5788	exact à gauche et accessible.
  5789	\end{enonce*}
  5790	
  5791	Comme\pageoriginale les foncteurs représentables
  5792	$h_{Y'}:X'\rightarrow \Hom(Y',X'):C'\rightarrow(\Ens)$ sont exacts à
  5793	gauche et accessibles (\hyperref[I.prop9.11]{9.11}), si $f$ a ces mêmes
  5794	propriétés, il en est de même de ses composés avec les foncteurs
  5795	précédents, qui sont donc proreprésentables par \hyperref[I.cor8.13.3]{8.13.3},
  5796	\ie $f$ admet un proadjoint. Inversement, supposons que $f$ admet un
  5797	proadjoint, et soient $C'_1$ une petite sous-catégorie génératrice
  5798	de $C$, $\Pi$ un cardinal tel que les $Y'\in \Ob C'$ soient
  5799	$\Pi$-accessibles; pour prouver que $f$ est $\Pi$-accessible, il
  5800	suffit donc de prouver qu'il en est ainsi de ses composés avec les
  5801	$h_{Y'}$, $Y'\in \Ob C'$; or par hypothèse ces composés sont
  5802	proreprésentables, donc ils sont accessibles (\hyperref[I.cor8.13.3]{8.13.3}),
  5803	donc $\Pi$-accessibles pourvu qu'on prenne $\Pi$ assez grand, \cqfd
  5804	
